% ============================================================
% 5.2 Hiện thực: Hoàn trả và Tranh chấp
% ============================================================

\section{Hiện thực --- Hoàn trả và Tranh chấp}
\label{sec:impl-refund}

\subsection{Mô tả cách sử dụng}

Buyer có thể yêu cầu hoàn trả cho đơn hàng đã thanh toán thành công. Quy trình gồm: buyer tạo yêu cầu hoàn trả (chọn phương thức PickUp hoặc DropOff, nhập lý do và ảnh minh chứng), seller xem xét và chấp nhận/từ chối. Nếu seller từ chối, buyer có quyền mở tranh chấp (dispute) để admin phân xử.

\begin{figure}[H]
\centering
\includegraphics[width=\textwidth,keepaspectratio]{uploads/refund.png}
\captionof{figure}{Form tạo yêu cầu hoàn trả}
\label{fig:impl-refund-form}
\end{figure}

Buyer điền form hoàn trả gồm phương thức trả hàng, lý do chi tiết và ảnh minh chứng. Hệ thống tự động gửi thông báo cho seller qua SSE.

\iffalse
\begin{center}
\fbox{\parbox{0.85\textwidth}{\centering\vspace{3cm}TODO: Screenshot trang tranh chấp --- hiển thị lý do từ chối của seller, form mở dispute với bằng chứng bổ sung\vspace{3cm}}}
\captionof{figure}{Trang tranh chấp hoàn trả (Refund Dispute)}
\label{fig:impl-refund-dispute}
\end{center}
\fi

Khi seller từ chối hoàn trả, buyer có thể mở tranh chấp bằng cách cung cấp bằng chứng bổ sung. Admin sẽ xem xét và đưa ra phán quyết chấp nhận hoặc từ chối.


% ------------------------------------------------------------------
\subsection{Đánh giá}
\label{subsec:refund-eval}

\small
\rowcolors{2}{white}{rowgray}
\begin{longtable}{@{}L{0.5cm}L{6.5cm}L{0.5cm}L{6.5cm}@{}}
\caption{Đánh giá chức năng hoàn trả và tranh chấp}\label{tab:refund-eval}\\
\toprule
\rowcolor{headerblue}
\multicolumn{2}{c}{\thead{Điểm mạnh}} & \multicolumn{2}{c}{\thead{Hạn chế}} \\
\midrule
\endfirsthead
\toprule
\rowcolor{headerblue}
\multicolumn{2}{c}{\thead{Điểm mạnh}} & \multicolumn{2}{c}{\thead{Hạn chế}} \\
\midrule
\endhead
\midrule
\multicolumn{2}{r}{\small\textit{Tiếp theo trang sau}} \\
\endfoot
\bottomrule
\endlastfoot

1 & Hỗ trợ hai phương thức trả hàng (PickUp/DropOff) linh hoạt & 1 & Chưa tích hợp hoàn tiền tự động qua cổng thanh toán \\
2 & Cơ chế dispute ba bên (buyer--seller--admin) công bằng & 2 & Chưa có thời hạn tự động duyệt nếu seller không phản hồi \\
3 & Upload ảnh minh chứng qua presigned URL (MinIO) & 3 & Chưa hỗ trợ hoàn trả một phần (partial refund) \\
4 & Thông báo real-time qua SSE cho các bên liên quan & 4 & Chưa có luồng trả hàng tích hợp với đơn vị vận chuyển \\
\end{longtable}
\rowcolors{1}{}{}
\normalsize

% ------------------------------------------------------------------
\subsection{Kiểm thử hộp đen}
\label{subsec:refund-blackbox}

Kiểm thử hộp đen cho hoàn trả và tranh chấp bao gồm hai nhóm: luồng hoàn trả (tạo, sửa, hủy, duyệt) và luồng tranh chấp (mở, phán quyết). Test case được thiết kế theo \textbf{phân lớp tương đương} kết hợp kiểm tra phân quyền (buyer / seller / admin / người ngoài).

\small
\rowcolors{2}{white}{rowgray}
\begin{longtable}{@{}L{1.2cm}L{3cm}L{4cm}L{4.5cm}L{0.8cm}@{}}
\caption{Test case hộp đen --- hoàn trả và tranh chấp}\label{tab:refund-test}\\
\toprule
\rowcolor{headerblue}
\thead{ID} & \thead{Kịch bản} & \thead{Điều kiện trước / Đầu vào} & \thead{Kết quả mong đợi} & \thead{KQ} \\
\midrule
\endfirsthead
\toprule
\rowcolor{headerblue}
\thead{ID} & \thead{Kịch bản} & \thead{Điều kiện trước / Đầu vào} & \thead{Kết quả mong đợi} & \thead{KQ} \\
\midrule
\endhead
\midrule
\multicolumn{5}{r}{\small\textit{Tiếp theo trang sau}} \\
\endfoot
\bottomrule
\endlastfoot
\multicolumn{5}{@{}l}{\textit{Hoàn trả}} \\
TC-15 & Tạo hoàn trả PickUp & Đơn hàng đã giao; method = PickUp, có địa chỉ & Yêu cầu tạo thành công, seller nhận thông báo & Đạt \\
TC-16 & Tạo hoàn trả PickUp thiếu địa chỉ & method = PickUp, address trống & Bị từ chối: ``Address is required'' & Đạt \\
TC-17 & Tạo hoàn trả DropOff & method = DropOff, không cần địa chỉ & Yêu cầu tạo thành công & Đạt \\
TC-18 & Hoàn trả một phần (partial) & Chọn 2/5 items, nhập số tiền cụ thể & Hoàn trả tạo với item\_ids và amount đúng & Đạt \\
TC-19 & Hoàn trả với item\_ids trùng lặp & item\_ids = [1, 1, 2] & Bị từ chối: ``Duplicate item IDs'' & Đạt \\
TC-20 & Sửa hoàn trả đang Pending & Đổi lý do và ảnh minh chứng & Cập nhật thành công, ảnh mới thay thế & Đạt \\
TC-21 & Sửa hoàn trả đã Processing & Hoàn trả đã được seller duyệt & Bị từ chối (409): không thể sửa & Đạt \\
TC-22 & Buyer hủy hoàn trả & Hoàn trả đang Pending & Chuyển Cancelled, seller nhận thông báo & Đạt \\
TC-23 & Seller chấp nhận hoàn trả & Hoàn trả Pending, seller nhấn ``Duyệt'' & Chuyển Processing, buyer nhận thông báo & Đạt \\
TC-24 & Seller từ chối hoàn trả & Hoàn trả Pending, seller nhấn ``Từ chối'' & Chuyển Cancelled, buyer nhận thông báo & Đạt \\
\midrule
\multicolumn{5}{@{}l}{\textit{Tranh chấp}} \\
TC-25 & Buyer mở dispute & Hoàn trả bị từ chối (Cancelled), có lý do & Dispute tạo, seller nhận thông báo & Đạt \\
TC-26 & Seller mở dispute & Hoàn trả đang Processing, seller bất đồng & Dispute tạo, buyer nhận thông báo & Đạt \\
TC-27 & Admin phán quyết chấp nhận & Dispute Pending, admin chọn ``Chấp nhận'' & Dispute resolved, hoàn trả tiếp tục xử lý & Đạt \\
TC-28 & Admin phán quyết từ chối & Dispute Pending, admin chọn ``Từ chối'' & Dispute đóng, hoàn trả bị từ chối & Đạt \\
TC-29 & Mở dispute khi đã có dispute đang hoạt động & Dispute Pending tồn tại cho hoàn trả này & Bị từ chối (409): ``Active dispute exists'' & Đạt \\
TC-30 & Người ngoài mở dispute & Account không phải buyer/seller của đơn & Bị từ chối (403): ``Not authorized'' & Đạt \\
\end{longtable}
\rowcolors{1}{}{}
\normalsize
